\input{configTD}

\title{Probabilités TD1 - Evénements et probabilités}

\author{}
\date{}

\begin{document}
\sffamily
\maketitle

\section*{1 - Manipulation d'événements}
On contrôle successivement une suite de pièces utilisées sur un chantier. Pour $n \geq 1$, on note $A_n$ l'événement "la $n$-ième pièce contrôlée est conforme".
\begin{enumerate}
    \item Comment écrire à partir des $A_i$ les événements :
    \begin{itemize}
        \item Les 3 premières pièces contrôlées sont conformes
        \item Parmi les 2 premières pièces contrôlées, la première est conforme et la seconde ne l'est pas.
        \item Parmi les 2 premières pièces contrôlées, 1 est conforme et 1 ne l'est pas.
    \end{itemize}

    \ifthenelse{\boolean{showSolutions}}{
        \begin{solution}
    \begin{itemize}
        \item $A_1 \cap A_2 \cap A_3$
        \item $A_1 \cap \bar{A_2}$
        \item $(A_1\cap \bar{A_2}) \cup (\bar{A_1}\cap A_2)$
    \end{itemize}
\end{solution}
    }{}
    \item Décrire par une phrase les événements suivants:
    $$
B=\bigcap_{i=5}^{+\infty} A_i,\qquad  C=\bigcup_{i=5}^{+\infty} A_i
$$
    \ifthenelse{\boolean{showSolutions}}{
        \begin{solution}
    $B$ signifie qu'à partir du cinquième contrôle, toutes les pièces contrôlées sont conformes. $C$ signifie qu'au moins une pièce contrôlée à partir du cinquième contrôle est conforme. On peut décrire $D$ par la formule
        \end{solution}
    }{}


\item  Écrire à l'aide de $A_i$ l'événement $D$ "À partir d'un certain contrôle, toutes les pièces contrôlées sont conformes".

\ifthenelse{\boolean{showSolutions}}{
    \begin{solution}
Si on avait voulu l'évennement $D_{10}$: "A partir du 10eme contrôle, toutes les pièces sont conformes", on écrirait 
$$
D_{10}=\bigcap_{i=10}^{+\infty} A_i
$$
Pour écrire $D$, il faut considérer que c'est :
\begin{enumerate}
    \item soit à partir du premier, toutes sont conformes,
    \item soit à partir du deuxième, toutes sont conformes, 
    \item etc .. 
    \item soit à partir du dixième toutes sont conformes, 
    \item etc ...
\end{enumerate}

C'est donc la réunion de tous ces évennements: 

$$
D=\bigcup_{n=1}^{+\infty} \bigcap_{i \geq n} A_i\,.
$$

    \end{solution}
}{}

\item On suppose que chaque pièce a une probabilité $\frac{1}{2}$ d'être conforme, indépendamment des autres. Calculer la probabilité que, parmi les deux premières pièces contrôlées, l'une au moins soit conforme.

\ifthenelse{\boolean{showSolutions}}{
    \begin{solution}
On note $A_1$ l'événement "la première pièce est conforme" et $A_2$ l'événement "la deuxième pièce est conforme". On cherche

$$
P(A_1 \cup A_2)=1-P(\bar{A_1}\cap \bar{A_2}).
$$

Par indépendance,

$$
P(\bar{A_1}\cap \bar{A_2})=P(\bar{A_1})P(\bar{A_2})=\frac{1}{2}\times \frac{1}{2}=\frac{1}{4}.
$$

Donc

$$
P(A_1 \cup A_2)=1-\frac{1}{4}=\frac{3}{4}.
$$
    \end{solution}
}{}

\item Quelle est la probabilité que la deuxième pièce soit conforme sachant que la première l'est ?

\ifthenelse{\boolean{showSolutions}}{
    \begin{solution}
Comme les contrôles sont indépendants,

$$
P(A_2 \mid A_1)=P(A_2)=\frac{1}{2}.
$$

Autrement dit, savoir que la première pièce est conforme ne change pas la probabilité que la deuxième le soit.
    \end{solution}
}{}

\item Sachant qu'une des deux premières pièces est conforme, quelle est la probabilité que l'autre le soit aussi ?
\ifthenelse{\boolean{showSolutions}}{
    \begin{solution}
L'information "une des deux premières pièces est conforme" correspond à l'événement $A_1 \cup A_2$. Dire que l'autre l'est aussi signifie alors que les deux sont conformes, c'est-à-dire l'événement $A_1 \cap A_2$.

On cherche donc

$$
P(A_1 \cap A_2 \mid A_1 \cup A_2)=\frac{P(A_1 \cap A_2)}{P(A_1 \cup A_2)}.
$$

Or, par indépendance,

$$
P(A_1 \cap A_2)=\frac{1}{2}\times \frac{1}{2}=\frac{1}{4}
$$

et, d'après la question précédente,

$$
P(A_1 \cup A_2)=\frac{3}{4}.
$$

Ainsi,

$$
P(A_1 \cap A_2 \mid A_1 \cup A_2)=\frac{\frac{1}{4}}{\frac{3}{4}}=\frac{1}{3}.
$$
    \end{solution}
}{}
\end{enumerate}
\vspace{2em}

\section*{2 - Mesure de Probabilité}
Soit $a \in] 0,1[$.

Sur $\mathbb{N}$, une mesure de probabilité est entièrement déterminée par les poids $P(\{k\})$ pour $k \in \mathbb{N}$, avec
\[
P(\{k\}) \geq 0
\qquad \text{et} \qquad
\sum_{k=0}^{+\infty} P(\{k\}) = 1.
\]
Pour toute partie $A \subset \mathbb{N}$, on a alors
\[
P(A)=\sum_{k \in A} P(\{k\}).
\]
\begin{enumerate}

\item Déterminer $a$ tel que la fonction $P$, qui à tout $n \in \mathbb{N}$ associe $P(\{n\})=(1-a) a^n$ soit une mesure de probabilité sur $\mathbb{N}$. 
\ifthenelse{\boolean{showSolutions}}{
    \begin{solution}
 Il suffit de vérifier que, pour tout $\left.n \in \mathbb{N}, a^n(1-a) \in\right] 0,1[$ - ce qui est clair puisque $a \in] 0,1[$, donc $\left.\left.a^n \in\right] 0,1\right]$ ( $n$ peut être égal à 0 ) et $\left.1-a \in\right] 0,1[$ - et que

$$
\sum_{n=0}^{+\infty}(1-a) a^n=\frac{1-a}{1-a}=1
$$
    \end{solution}
}{}

\item On considère les deux événements $A=\{2 k: k \in \mathbb{N}\}$ et $B=\{2 k+1: k \in \mathbb{N}\}$. Calculer $P(A)$ et $P(B)$. Les événements $A$ et $B$ sont-ils incompatibles ? sont-ils indépendants ?
\end{enumerate} 

\ifthenelse{\boolean{showSolutions}}{
    \begin{solution}
On a 

$$
P(A)=\sum_{k=0}^{+\infty}(1-a) a^{2 k}=\frac{1-a}{1-a^2}=\frac{1}{1+a}
$$


Comme $B=\bar{A}$, on a

$$
P(B)=1-P(A)=\frac{a}{1+a}
$$


Puisque $A \cap B=\varnothing$, les événements $A$ et $B$ sont incompatibles. Puisque $P(A \cap B)=P(\varnothing)=0$ et $P(A) P(B) \neq 0$, les événements $A$ et $B$ ne sont pas indépendants.
    \end{solution}
}{}
\vspace{2em}

\section*{3 - Indépendance d'événements}
\begin{enumerate}
\item Un atelier fabrique $12$ pièces numérotées de $1$ à $12$. On en prélève une au hasard pour contrôle, et on considère les événements

$$
\begin{aligned}
& A=\text { "la pièce prélevée porte un numéro pair", } \\
& B=\text { "la pièce prélevée porte un numéro multiple de } 3 \text { ". }
\end{aligned}
$$


Les événements $A$ et $B$ sont-ils indépendants?

\ifthenelse{\boolean{showSolutions}}{
    \begin{solution}
        On a :

$$
\begin{gathered}
A=\{2,4,6,8,10,12\} \\
B=\{3,6,9,12\} \\
A \cap B=\{6,12\}
\end{gathered}
$$
On a donc $P(A)=1 / 2, P(B)=1 / 3$ et $P(A \cap B)=1 / 6=P(A) P(B)$. Les événements $A$ et $B$ sont indépendants.

\end{solution}
}{}

\item  Reprendre la question avec un atelier fabriquant $13$ pièces numérotées.
\end{enumerate} 

\ifthenelse{\boolean{showSolutions}}{
    \begin{solution}
Les événements $A, B$ et $A \cap B$ s'écrivent encore exactement de la même façon. Mais cette fois, on a : $P(A)=6 / 13, P(B)=4 / 13$ et $P(A \cap B)=2 / 13 \neq 24 / 169$. Les événements $A$ et $B$ ne sont pas indépendants. C'est conforme à l'intuition. It n'y a plus la même répartition de boules paires et de boules impaires, et dans les multiples de 3 compris entre 1 et 13 , la répartition des nombres pairs et impairs est restée inchangée.
    \end{solution}
}{}
\vspace{2em}

\section*{4 - Lois discrètes classiques}
Dans cet exercice, on cherche \`a reconna\^itre dans quelles situations les lois discr\`etes classiques fournissent un mod\`ele adapt\'e.

\begin{enumerate}
    \item Relier chaque loi discr\`ete de la colonne de gauche \`a la situation correspondante de la colonne de droite. On justifiera bri\`evement chaque choix.
    
    \medskip
    \noindent
    \begin{minipage}[t]{0.33\textwidth}
    \textbf{Colonne A -- Lois}\vspace{1em}
    
    Loi binomiale $B(n, p)$ \hfill \rule{1.2ex}{1.2ex}
    
    \medskip
    Loi de Poisson $P(\lambda)$ \hfill \rule{1.2ex}{1.2ex}
    
    \medskip
    Loi g\'eom\'etrique $G(p)$ \hfill \rule{1.2ex}{1.2ex}
    
    \medskip
    Loi uniforme discr\`ete $U(a, b)$ \hfill \rule{1.2ex}{1.2ex}
    \end{minipage}
    \hspace{0.14\textwidth}
    \begin{minipage}[t]{0.43\textwidth}
    \textbf{Colonne B -- Situations}\vspace{1em}
    
    \rule{1.2ex}{1.2ex}\quad Nombre de reprises parmi $20$ interventions.
    
    \medskip
    \rule{1.2ex}{1.2ex}\quad Nombre d'incidents rares en une journ\'ee.
    
    \medskip
    \rule{1.2ex}{1.2ex}\quad Rang du premier contr\^ole concluant.
    
    \medskip
    \rule{1.2ex}{1.2ex}\quad Choix al\'eatoire d'un cr\'eneau entre $t_1$ et $t_2$.
    \end{minipage}
    
    \vspace{2em}

    \item Pour chacune des variables al\'eatoires suivantes, dire quelle loi discr\`ete permet de la mod\'eliser et pr\'eciser ses param\`etres.
    \begin{enumerate}
        \item $X$ d\'esigne le nombre d'\'el\'ements non conformes parmi $12$ \'el\'ements contr\^ol\'es, chaque \'el\'ement ayant une probabilit\'e $0{,}03$ d'\^etre non conforme ;
        \item $Y$ d\'esigne le nombre de jours de pluie pendant une semaine de $7$ jours, chaque jour ayant une probabilit\'e $0{,}3$ d'\^etre pluvieux, ind\'ependamment des autres ;
        \item $Z$ d\'esigne le nombre d'incidents mineurs observ\'es pendant une semaine sur un chantier, sachant qu'on en observe en moyenne $1{,}2$ par semaine ;
        \item $T$ d\'esigne le rang du premier essai concluant lors d'une s\'erie d'essais ind\'ependants, chaque essai ayant une probabilit\'e $0{,}75$ de r\'eussir ;
        \item $U$ d\'esigne le num\'ero du quai de d\'echargement attribu\'e au hasard \`a un camion parmi les quais num\'erot\'es de $1$ \`a $6$ ;
        \item $V$ d\'esigne le nombre d'appels de maintenance re\c cus en une journ\'ee, sachant que ce nombre vaut en moyenne $3$ par jour.
    \end{enumerate}
    
\end{enumerate}

\ifthenelse{\boolean{showSolutions}}
{
    \begin{solution}
        \begin{enumerate}
            \item Correspondances :
            \[
            1 \leftrightarrow A,
            \qquad
            2 \leftrightarrow B,
            \qquad
            3 \leftrightarrow C,
            \qquad
            4 \leftrightarrow D.
            \]
            En effet, $A$ correspond \`a un nombre de succ\`es dans un nombre fix\'e d'essais ind\'ependants, $B$ \`a un nombre d'\'ev\'enements rares sur une dur\'ee donn\'ee, $C$ au rang du premier succ\`es, et $D$ \`a un choix \'equiprobable parmi des entiers.
            
            \item Mod\'elisation des variables al\'eatoires :
            \begin{enumerate}
                \item $X \sim B(12,0{,}03)$ ;
                \item $Y \sim B(7,0{,}3)$ ;
                \item $Z \sim P(1{,}2)$ ;
                \item $T \sim G(0{,}75)$ ;
                \item $U \sim U(1,6)$ ;
                \item $V \sim P(3)$.
            \end{enumerate}
            
        \end{enumerate}
    \end{solution}
}{}

\newpage

\section*{5 - Rebuts}
Une fabrication automatique de pièces embouties donne un pourcentage de rebuts s'élevant à $5 \%$. On considère un échantillon de $10$ pièces issues de cette fabrication. Calculer la probabilité de trouver dans cet échantillon au plus $2$ rebuts.

\ifthenelse{\boolean{showSolutions}}{
    \begin{solution}
Soit $X$ le nombre de rebuts parmi les $10$ pièces. Chaque pièce est défectueuse indépendamment avec probabilité $0{,}05$, donc $X \sim B(10;\, 0{,}05)$.

On cherche $P(X \leq 2) = P(X=0)+P(X=1)+P(X=2)$. Par la formule binomiale :
\begin{align*}
P(X=0) &= 0{,}95^{10} \approx 0{,}5987\,,\\
P(X=1) &= \binom{10}{1}\times 0{,}05 \times 0{,}95^{9} \approx 0{,}3151\,,\\
P(X=2) &= \binom{10}{2}\times 0{,}05^{2} \times 0{,}95^{8} \approx 0{,}0746\,.
\end{align*}
Donc $P(X \leq 2) \approx 0{,}9884$.
    \end{solution}
}{}

\vspace{1em}

\section*{6 - Un dé}

Combien de fois faut-il lancer un dé pour faire au moins un six avec une probabilité supérieure ou égale à $0{,}95$ ?

\ifthenelse{\boolean{showSolutions}}{
    \begin{solution}
L'événement contraire est "aucun six en $n$ lancers", de probabilité $\left(\tfrac{5}{6}\right)^n$. On cherche le plus petit entier $n$ tel que
$$
1 - \left(\frac{5}{6}\right)^n \geq 0{,}95
\iff
\left(\frac{5}{6}\right)^n \leq 0{,}05.
$$
En passant au logarithme (en base $e$) :
$$
n \geq \frac{\ln(0{,}05)}{\ln\!\left(\tfrac{5}{6}\right)} = \frac{\ln(0{,}05)}{\ln 5 - \ln 6} \approx \frac{-2{,}996}{-0{,}182} \approx 16{,}4\,.
$$
Il faut donc au moins $\mathbf{17}$ lancers.
    \end{solution}
}{}

\vspace{1em}

\section*{7 - Formule des probabilités totales}
Dans une entreprise deux ateliers fabriquent les mêmes pièces. 
L'atelier 1 fabrique en une journée deux fois plus de pièces que l'atelier 2. 
Le pourcentage de pièces défectueuses est $3 \%$ pour l'atelier 1 et $4 \%$ pour l'atelier 2. 
On prélève une pièce au hasard dans l'ensemble de la production d'une journée. 

Déterminer (\textit{On pourra représenter un arbre de probabilités})
\begin{enumerate}
    \item la probabilité que cette pièce provienne de l'atelier 1 ;
    \ifthenelse{\boolean{showSolutions}}{
        \begin{solution}
            Notons $A$ l'événement " la pièce provient de l'atelier 1 ", $B$ l'événement " la pièce provient de l'atelier 2" et $D$ l'événement " la pièce est défectueuse".
    
            L'énoncé nous dit que les $2 / 3$ des pièces produites proviennent de l'atelier 1. Donc $P(A)=2 / 3$
        \end{solution}
    }{}

    \item  la probabilité que cette pièce provienne de l'atelier 1 et est défectueuse;
    \ifthenelse{\boolean{showSolutions}}{
        \begin{solution}
    On cherche $P(A \cap D)=P_A(D) P(A)=0,03 \times \frac{2}{3}=\frac{1}{50}$.
        \end{solution}
    }{}

    \item  la probabilité que cette pièce provienne de l'atelier 1 sachant qu'elle est défectueuse.
\end{enumerate}

\ifthenelse{\boolean{showSolutions}}{
    \begin{solution}
On démontre de même que $P(B \cap D)=\frac{1}{75}$ et donc que

$$
P(D)=P(A \cap D)+P(B \cap D)=\frac{1}{30}
$$


Ainsi, on a

$$
P_D(A)=\frac{P(A \cap D)}{P(D)}=\frac{3}{5}
$$

    \end{solution}
}{}
\vspace{1em}

\section*{8 - Loi de variables aléatoires}
Soit $X$ une variable aléatoire suivant une loi uniforme sur $\{1, \ldots, 20\}$. Déterminer la loi de $\lfloor\sqrt{X}\rfloor$.

Vous pourrez vérifier vos calculs en montrant que la loi de $\sqrt{X}$ est bien une mesure de probabilité, sur quel ensemble ?

Calculer l'espérance de $X$ et l'espérance de $\lfloor\sqrt{X}\rfloor$.

On donne la formule pour calculer l'espérance : si une variable aléatoire $X$ peut prendre les valeurs $x_1, x_2, \cdots$, alors
\[
E(X) =  \sum_{k\geq 1}x_k P(X=x_k)\,. 
\]
En français, l'espérance est la moyenne des valeurs que peut prendre $X$, pondérée par les probabilités que $X$ prenne ces valeurs. 

\ifthenelse{\boolean{showSolutions}}{
    \begin{solution}
        Pour simplifier les notations, posons $Y=\lfloor\sqrt{X}\rfloor$. Alors, $Y$ peut prendre les valeurs $1,2,3,4$. Plus précisément,

    \begin{itemize}

\item  si $X=1, X=2$ ou $X=3$, alors $Y=1$. En particulier,

$$
P(Y=1)=P(X=1)+P(X=2)+P(X=3)=\frac{3}{20}
$$

\item  si $X=4, X=5, X=6, X=7$ ou $X=8$, alors $Y=2$. En particulier,

$$
P(Y=2)=\frac{5}{20}
$$

\item si $X=9, X=10, X=11, X=12, X=13, X=14$ ou $X=15$, alors $Y=3$. En particulier,

$$
P(Y=3)=\frac{7}{20}
$$

\item si $X=16,17,18,19$ ou 20 , alors $Y=4$. En particulier,

$$
P(Y=4)=\frac{5}{20}
$$
    \end{itemize}

Pour calculer l'espérance, on utilise la formule : $X$ peut prendre toutes les valeurs entre $1$ et $20$, chacune avec probabilité $1/20$, l'espérance de $X$ est donc : 
\begin{align*}
E(X) &=  1 \times  \frac{1}{20} +  2 \times  \frac{1}{20}  +  3 \times  \frac{1}{20}  + \cdots +  20 \times  \frac{1}{20} \\
    & = \frac{1}{20} \sum_{k=1}^{20} k \\
    & = \frac{1}{20} \frac{20\times 21}{2} \\ 
    & = \frac{21}{2}\,.
\end{align*}

Calculons maintenant l'espérance de $Y$ : on multiplie les valeurs que peut prendre $Y$ par les probabilités qu'on a calculé : 
\[
E(Y) = 1 \times \frac{3}{20} + 2 \times \frac{5}{20} + 3\times \frac{7}{20} + 4 \times \frac{5}{20}
\]


    \end{solution}
}{}

\newpage

\section*{9 - Un avant-goût de statistiques}
Un étang contient des brochets et des truites. On note $p$ la proportion de truites dans l'étang. On souhaite évaluer $p$, on prélève 20 poissons au hasard. On suppose que le nombre de poissons est suffisamment grand pour que ce prélèvement s'apparente à 20 tirages indépendants avec remise. 
On note $X$ le nombre de truites obtenues.
\begin{enumerate}
\item Quelle est la loi de $X$ ?
\item Le prélèvement a donné 8 truites. Pour quelle valeur de $p$ la quantité $P(X=8)$ est-elle maximale?
\item On réalise un deuxième prélèvement de $20$ poissons qui a donné $9$ truites, comment modifier votre estimation pour $p$ ?
\item On réalise $n$ prélèvements de $20$ poissons, comment obtenir une bonne estimation de $p$ ?
\end{enumerate}
\ifthenelse{\boolean{showSolutions}}{
    \begin{solution}
        1. On reconnait le principe d'un schéma de Bernoulli, avec 20 épreuves indépendantes et une probabilité $p$ de succès. $X$ suit donc une loi binomiale $\mathcal{B}(20, p)$.

        2. On a $P(X=8)=\binom{20}{8} p^8(1-p)^{12}$. On va étudier cette fonction de $p$ de sorte de trouver son maximum. Il suffit en réalité d'étudier, sur $[0,1]$, la fonction $g(p)=p^8(1-p)^{12} . g$ est dérivable et sa dérivée est

$$
\begin{aligned}
g^{\prime}(p) & =8 p^7(1-p)^{12}-12 p^8(1-p)^{11} \\
& =(8(1-p)-12 p) p^7(1-p)^{11} \\
& =(8-20 p) p^7(1-p)^{11}
\end{aligned}
$$


On en déduit que $g^{\prime}(p) \geq 0$ sur $[0,2 / 5]$ et $g^{\prime}(p) \leq 0$ sur $[2 / 5,1]$. La probabilité $P(X=8)$ est maximale pour $p=2 / 5$.
    \end{solution}
}{}
\vspace{2em}


\end{document}